\begin{proposicion}[Cociente de normalizadores]\label{coste-normalizadores}
Bajo las mismas hipótesis se tiene
\[
 \log(\Gamma_4/\Gamma_2)=O_{C,\eta,\gamma}(\varepsilon_{\rm out}+k^{-1}).
\]
\end{proposicion}
\begin{proof}
Escribimos $a_2=\sum_e p_e^2$. Primero comparamos los bloques de filas y columnas. La derivada de
$\log\det$ en el complemento
ortogonal del núcleo común es la traza del producto de la inversa por la derivada de la matriz. El
Lema~\ref{coste-inversa} y $\tau_4(p)-\tau_2(p)=2p^2+O(p^3)$ dan
\[
 \begin{split}
 \log\frac{\pdet S_{R,4}}{\pdet S_{R,2}}
 &=\int_0^1\sum_e(\tau_4(p_e)-\tau_2(p_e))\,
                   \iota_e^TS(s)^+\iota_e\,ds\\
 &=4a_2+O(k^{-1}).
 \end{split}
\]
En efecto, $\iota_e^TS(s)^+\iota_e=2+O(k^{-1})$,
$a_2=O(1)$ y $\sum_e p_e^3\le p_{\max}a_2=O(k^{-1})$.

Escribamos $\tau_e=\tau_4(p_e)$, $G=S_{R,4}^+$ y
\[
 V=A_R^T\operatorname{diag}(\tau)D,\qquad
 S_D=\operatorname{diag}(h)+D^T\operatorname{diag}(\tau)D,\qquad
 U=S_D-V^TGV .
\]
Cada $\iota_e$ es ortogonal al modo cero y las columnas de $V$ son combinaciones de esos vectores. En
una base ortonormal del complemento del modo cero, la eliminación por bloques da
$\pdet S_4=(\pdet S_{R,4})\det U$. Para cualquier $z$ diagonal,
la forma $z^TU z$ es el mínimo de la forma de $S_4$ en $(y,z)$ sobre $y$ ortogonal al modo cero. El
minimizador es $-GVz$. La cota inferior del Lema~\ref{coste-estatico} implica por tanto $U\ge cI$ con
el mismo $c>0$.

Pongamos $A=S_D-I$ y $T=V^TGV$, de modo que $U=I+A-T$. En la diagonal,
$A_{ff}=\sum_{e\in f}(\tau_e-p_e)$. Fuera de ella sólo hay entradas entre las dos familias diagonales,
con valor $\tau_e$ cuando las líneas se cortan en la casilla disponible $e$ y cero en otro caso. Se
sigue que
\[
 \|A\|_{\rm op}=O(\varepsilon_{\rm out}),\qquad
 \operatorname{tr}A=6a_2+O(k^{-1}),\qquad
 \operatorname{tr}A^2=2a_2+O(k^{-1}).
\]
En efecto, la suma absoluta de cada fila es $O(\ell_f)$, y la traza cuenta dos veces
$\tau_e-p_e=3p_e^2+O(p_e^3)$ por casilla. Para el cuadrado, las entradas cruzadas aportan
$2\sum_e\tau_e^2=2a_2+O(k^{-1})$. La contribución de la diagonal es $O(k^{-1})$ por $\max_f
A_{ff}=O(k^{-1})$ y $\sum_f A_{ff}=O(1)$.

Cada entrada no nula de $V$ es $\tau_e$, ya que una fila o columna y una diagonal determinan a lo sumo
una casilla. Las sumas en las filas de $V$ son a lo sumo $2\max_u\sum_{e\in u}\tau_e=O(1)$, donde $u$
recorre filas y columnas. Las sumas en columnas son a lo sumo $2\max_f\sum_{e\in
f}\tau_e=O(\varepsilon_{\rm out})$. Por la desigualdad entre la norma espectral y las normas de sumas
de filas y columnas,
$\|V\|_{\rm op}^2=O(\varepsilon_{\rm out})$. Además
$\|V\|_{\rm HS}^2=4\sum_e\tau_e^2=O(1)$. Como $G$ es positivo semidefinido y tiene norma a lo sumo
$\gamma^{-1}$,
\[
 T\ge0,\qquad \|T\|_{\rm op}=O(\varepsilon_{\rm out}),\qquad
 \operatorname{tr}T=O(1),\qquad
 \operatorname{tr}T^2\le\|T\|_{\rm op}\operatorname{tr}T
             =O(\varepsilon_{\rm out}) .
\]

El término principal de la traza de $T$ requiere usar la incidencia. Por el Lema~\ref{coste-inversa},
\[
 \operatorname{tr}T
 =\sum_f\sum_{e,e'\in f}\tau_e\tau_{e'}\iota_e^TG\iota_{e'}
 =4a_2+O(\varepsilon_{\rm out}+k^{-1}) .
\]
Los pares $e=e'$ aportan $2\sum_e\tau_e^2(2+O(k^{-1}))$: hay dos diagonales por casilla. Dos casillas
distintas de la misma diagonal tienen filas y columnas distintas. Su forma mixta es por ello
$O(k^{-1})$. El total absoluto de estos errores se acota por
\[
 \frac{C'}k\sum_f\Bigl(\sum_{e\in f}\tau_e\Bigr)^2
 \le\frac{C''}k\sum_f\ell_f^2
 \le\frac{C''}k\,\varepsilon_{\rm out}\sum_f\ell_f
 =O(\varepsilon_{\rm out}),
\]
pues $\sum_f\ell_f=2k$.

Para $E=A-T$ tenemos $\|E\|_{\rm op}=O(\varepsilon_{\rm out})$ y $\operatorname{tr}E^2=O(1)$. También
$|\operatorname{tr}(AT)|\le\|A\|_{\rm op}\operatorname{tr}T=O(\varepsilon_{\rm out})$ por la
positividad de $T$. Por consiguiente,
\[
 \operatorname{tr}E=2a_2+O(\varepsilon_{\rm out}+k^{-1}),\qquad
 \operatorname{tr}E^2=2a_2+O(\varepsilon_{\rm out}+k^{-1}).
\]
Si $\xi$ es un autovalor de $E$, entonces $1+\xi\ge c$. Taylor da
$|\log(1+\xi)-\xi+\xi^2/2|\le|\xi|^3/(3\min\{c,1\}^3)$.
Al sumar, el error es $O(\|E\|_{\rm op}\operatorname{tr}E^2)=O(\varepsilon_{\rm out})$. En particular
\[
 \log\det U=a_2+O(\varepsilon_{\rm out}+k^{-1}).
\]
Por último, el coeficiente cuadrático de $\log(1-dp)-d\log(1-p)$ es $-d(d-1)/2$ y su resto es $O(p^3)$
para $d=2,4$. Así $C_4-C_2=-5a_2+O(k^{-1})$. El factor $\sqrt{4\pi k}$ es común a $\Gamma_4$ y
$\Gamma_2$. Reuniendo los tres términos calculados,
\[
 \log(\Gamma_4/\Gamma_2)
 =-\tfrac12\bigl((-5+4+1)a_2\bigr)
      +O(\varepsilon_{\rm out}+k^{-1}),
\]
como se afirmaba.
\end{proof}
